\chapter{Suíte de testes e evidências}
\label{ap:testes}

Este apêndice descreve a suíte automatizada e os instrumentos de evidência
utilizados no trabalho.

\section*{Cenas de teste do motor (uma por etapa)}
\label{ap:cenas}

Cada cena de teste instancia o jogo, executa interações reais e verifica o
estado resultante, sem interface gráfica e sem intervenção manual. A
execução completa é feita por um único comando
(\texttt{tests/run\_all.sh}), que percorre as cenas em ordem e, ao final,
executa também os testes do adaptador de IA.

\begin{itemize}
  \item \textbf{\poc{1} (77 verificações)} --- entrada semântica, movimento,
  colisão, gravidade, câmera, pausa, HUD, zona morta e tratamento de
  \emph{hotplug};
  \item \textbf{\poc{2} (34)} --- interação por proximidade, inventário,
  coleta, instalação no painel, avanço da missão, estados de equipamento,
  porta (incluindo a verificação de regressão da folha visual), iluminação em
  camadas, letreiros e câmera de segurança;
  \item \textbf{\poc{3} (25--26)} --- disponibilidade do serviço, limite de
  tempo, modelo ausente, resposta inválida, JSON malformado, modelo lento e
  servidor desligado, com verificação do \emph{fallback};
  \item \textbf{\poc{4} (21)} --- fala do NPC por estado, luzes conforme
  energia, tom da IA após alerta ignorado, quatro camadas de memória e limite
  do histórico;
  \item \textbf{\poc{5} (23)} --- módulos, pontos de interesse, determinismo
  por semente (seis sementes) e conectividade por busca em largura;
  \item \textbf{\poc{6} (30)} --- tutorial, missão completa por interação
  real, áudio (efeitos e ambiência em laço completo), salvamento e
  carregamento, geração da ala procedural, núcleo da IA e escolhas do
  epílogo.
\end{itemize}

\section*{Testes do adaptador de IA}
\label{ap:adapter}

Dezesseis testes em Python exercitam o serviço local: contrato de entrada e
saída, limites de tamanho da mensagem, validação de dados, conversão de
intenção, resposta de erro do modelo, servidor indisponível e formato de
resposta de \emph{fallback}.

\section*{Arnês de evidências visuais}
\label{ap:arnes}

Uma cena dedicada posiciona o jogador em coordenadas pré-determinadas,
opcionalmente conclui a missão (energia restabelecida e porta aberta) e salva
imagens em 1280$\times$720. Variáveis de ambiente controlam o caminho do
arquivo, a pose do jogador (posição e orientação) e a conclusão automática da
missão; há ainda modos auxiliares para inspecionar o terminal da IA e para
listar os elementos do cenário (NPC, luzes, letreiros e telas). As figuras
dos Capítulos~\ref{cap:desenvolvimento} e~\ref{cap:resultados} foram geradas
por esse instrumento.

\section*{Evidências por classe}
\label{ap:evidencias}

As evidências do trabalho se organizam em cinco classes: (i) testes
automatizados do motor; (ii) testes de contrato do adaptador; (iii) medições
de desempenho e latência; (iv) capturas de tela em poses controladas; e (v)
vídeo de demonstração. A matriz de rastreabilidade do projeto associa cada
afirmação do texto final a uma dessas classes (ver
MONOGRAFIA\_RASTREABILIDADE.md).
